\documentclass[11pt,compress,t,notes=noshow, xcolor=table]{beamer}

\input{../../style/preamble}
\input{../../latex-math/basic-math}
\input{../../latex-math/basic-ml}
\input{../../latex-math/optim-basics}

\title{Optimization in Machine Learning}

\begin{document}

\titlemeta{
Second order methods
}{
Limitations of Newton-Raphson
}{
figure/NR_saddle_2d.png
}{
\item Singular and indefinite Hessian
\item Non-descent directions
\item Saddle points
\item Computational cost
\vspace{0.5cm}
\item[] \hspace{-0.7cm}Slides based on\\
\hspace{-0.7cm}\furtherreading{AGGARWAL2020} (Ch. 5.6)
}

\begin{framei}[sep=L]{Limitation 1: Singular and indefinite Hessian}
\item \textbf{Problem}: Newton-Raphson is designed for \textit{convex quadratic} $f$ with pd Hessian, but
$\H = \nabla^2 f(\xv^{[t]})$ may be
\begin{itemizeM}
\item \textbf{singular}: $\nabla^2 f(\xv^{[t]}) \mathbf{d}^{[t]} = - \nabla f(\xv^{[t]})$ has no unique solution
$\Rightarrow$ step undefined
\item \textbf{near-singular}: step length blows up as $\lammin \to 0$, and errors already present in
$\H$ and $\nabla f$ (minibatch noise, finite-difference approximations, floating-point roundoff)
are amplified by $\kappa(\H) = |\lammax| / |\lammin|$
\item \textbf{indefinite}: $\mathbf{d}^{[t]}$ is not a descent direction (Limitation 2), saddle points (Limitation 3)
\end{itemizeM}
\item \textbf{When does (near-)singularity happen in ML?} Rank-deficient design, collinear features, $d > n$,
or degenerate optima
\end{framei}


\begin{framei}[sep=L]{Example: Singular Hessian in the $L_2$-SVM}
\item Unregularized $L_2$-SVM (squared hinge loss) with $\thetav \in \R^d$:
$$
\risket = \sumin \max\{1 - \yi \thetav^T \xi, ~0\}^2
$$
\item Only margin violators $\mathcal{M} = \{i: \yi \thetav^T \xi < 1\}$ contribute to $\H \in \R^{d \times d}$:
$$
\H = 2 \sum_{i \in \mathcal{M}} \xi (\xi)^T
$$
\item Each $\xi (\xi)^T$ has \textbf{rank 1} $\Rightarrow$ $\operatorname{rank}(\H) \leq |\mathcal{M}|$,
so we need at least $d$ margin violators for $\H$ to be invertible
\item \textbf{But}: the better our fit, the fewer violators
$\Rightarrow$ $\H$ becomes singular exactly where we want to converge
\end{framei}


\begin{framei}[sep=L]{Remedy: Regularized Newton step}
\item Replace $\H$ by $\H + \lambda \id$ with $\lambda > 0$: all eigenvalues are shifted by $\lambda$
$$
(\H + \lambda \id) \mathbf{d}^{[t]} = - \nabla f(\xv^{[t]})
$$
\item Choosing $\lambda$ slightly larger than $|\lammin|$ makes $\H + \lambda \id$ pd
$\Rightarrow$ invertible and descent direction
\item Alternative: use pseudoinverse $\H^+$ instead of $\H^{-1}$
\item \textbf{Caveat}: regularization does not cure ill-conditioning -- for nearly singular $\H$ the step
stays sensitive to errors in $\H$ and $\nabla f$
\end{framei}


\begin{framei}[sep=L]{Limitation 2: Non-descent directions}
\item \textbf{Problem}: In general, $\mathbf{d}^{[t]}$ is not a descent direction
\item \textbf{But}: If Hessian is positive definite, $\mathbf{d}^{[t]}$ is descent direction:
$$
\nabla f(\xv^{[t]})^T \mathbf{d}^{[t]} = - \nabla f(\xv^{[t]})^T (\nabla^2 f(\xv^{[t]}))^{-1} \nabla f(\xv^{[t]}) < 0
$$
\item Near minimum, Hessian is positive definite
\item For initial steps, Hessian is often not positive definite and Newton-Raphson may give non-descending update directions
\end{framei}


\begin{framei}[fs=small]{Example: Non-descent directions}
\item Consider $f(x) = \frac{x^4}{4} - x^2 + 2x$ with $f'(x) = x^3 - 2x + 2$, $f''(x) = 3x^2 - 2$
\item Starting at $x^{[0]} = 0$: $f''(0) = -2 < 0$ ($\H$ not pd) $\Rightarrow$ Newton-Raphson performs \textbf{ascent} step:
$$
x^{[1]} = x^{[0]} - \frac{f'(x^{[0]})}{f''(x^{[0]})} = 0 - \frac{2}{-2} = 1, \qquad f(x^{[1]}) = 1.25 > 0 = f(x^{[0]})
$$
\item From $x^{[1]} = 1$: $f''(1) = 1 > 0$, but the step leads right back to $x^{[2]} = 0$ \\
$\Rightarrow$ Newton-Raphson cycles forever between $0$ and $1$
\item \textbf{GD} from the same $x^{[0]} = 0$ only uses the neg. gradient $-f'(x^{[0]})=-2$ (descent direction) and converges for appropriate step size $\alpha$
\imageC[0.65]{figure/NR_divergence.png}
\end{framei}


\begin{framei}{Limitation 3: Saddle points}
\item \textbf{Problem}: Newton-Raphson solves $\nabla f(\xv) = \zero$ $\Rightarrow$ attracted to \textbf{any}
critical point (minima, maxima, saddles)
\item \textbf{Example}: $f(x_1, x_2) = x_1^2 - x_2^2$ with $\H = \left(\begin{smallmatrix} 2 & 0 \\ 0 & -2 \end{smallmatrix}\right)$
indefinite
\begin{itemizeS}
\item Since $f$ is quadratic, a single Newton step lands \textit{exactly} on the saddle $(0,0)$, from any starting point
\item GD is repelled along the negative curvature direction $x_2$ and escapes (as long as $x_2^{[0]} \neq 0$)
\end{itemizeS}
\imageC[0.72]{figure/NR_saddle_2d.png}
\end{framei}

\comment[TODO]{TH}{I am a little bit unhappy with this slide / the saddle point limitation in general (took it from Aggarwal but I think it is not very good). Missing a clear message particularly for the flip slide. Discuss and revise.}
\begin{framei}{Saddle points: Approximation flips}
\item Near an inflection point, the quadratic approximation flips, e.g., for $f(x) = x^3$:
\imageC[0.66]{figure/NR_saddle_1d.png}
\item At $x = 0$: $f'= f''= 0$ $\Rightarrow$ update $0/0$ undefined (\textbf{degenerate critical point})
\end{framei}


\begin{framei}{Saddle points: Why this matters}
\item In high-dimensional non-convex problems (e.g., DNNs), saddle points vastly outnumber true optima
\item \textbf{GD can escape} saddle points: it is not attracted by them but may still get stuck in flat regions
\item \textbf{Newton-Raphson cannot escape} saddle points: it is attracted \textit{indiscriminately} to all critical points
\item Second order is not uniformly better: it pays off for complex curvature, but not for
landscapes riddled with saddles
\item Practitioners often prefer first order methods with adaptive schemes (e.g., Adam), which pick up
some curvature information implicitly and at much lower cost (more on this in the next chapter)
\end{framei}


\begin{framei}{Limitation 4: Computational cost}
\item \textbf{Problem}: The Newton-Raphson update can be \textbf{computationally expensive}, as $\mathbf{d}^{[t]}$ is solution of the linear system
$$
\nabla^2 f(\xv^{[t]}) \mathbf{d}^{[t]} = - \nabla f(\xv^{[t]})
$$
\item For $\xv \in \R^d$, the Hessian has $d^2$ entries: $\mathcal{O}(d^2)$ memory, and $\mathcal{O}(d^3)$ per iteration to solve the linear system (compared to $\mathcal{O}(d)$ memory and no linear solve for GD)
\item For large $d$ (e.g., deep neural networks with $d \sim 10^9$ parameters),
      this is prohibitive: storing a dense $d \times d$ Hessian alone is infeasible
\item \textbf{In ML}: for ERM with $n$ observations, already forming $\H$ costs another $\mathcal{O}(n d^2)$ (compared to $\mathcal{O}(nd)$ for $\nabla \risket$)
\item Newton-Raphson needs far fewer iterations than GD, but for large $d$ the cost
\textit{per} iteration dominates
\item \textbf{Aim:} Find quasi-second order methods not relying on exact $\H$
\begin{itemizeS}
\item Quasi-Newton method
\item Gauss-Newton algorithm (for least squares)
\end{itemizeS}
\end{framei}

\endlecture
\end{document}
